\section{Pilha}

Uma pilha é uma estrutura de dados abstrata que define uma política na ordem das inserções
e remoções.
Essa política é denominada \textbf{LIFO}: o último elemento a entrar é sempre o primeiro
a ser removido.
Uma analogia é uma pilha de pratos num restaurante; normalmente, os pratos colocados
acima, que correspondem aos últimos a terem sido adicionados, são os primeiros a serem
removidos.
Ao manipular uma pilha, a operação de inserção é denominada \textit{push}, enquanto
a de retirada é denominada \textit{pop}.

Há pelo menos duas maneiras de implementar pilhas: com vetores ou com listas encadeadas
simples.
Independentemente da estrutura subjacente escolhida, é importante garantir que tanto a
adição quanto a remoção possuam custo constante em tempo.

Na implementação com vetor, basta adicionar o novo elemento ao final.
Dependendo do requisito da aplicação, realocações podem ou não ser realizadas.
Caso não o sejam, um erro de \textit{stack overflow} deve ser levantado.
Já para a operação de \textit{pop}, basta retirar o elemento na última posição.
Usualmente, seu conteúdo é retornado para impedir que ele seja perdido.
Dessa forma, cabe a quem chama o procedimento decidir se irá ou não salvar o objeto
devolvido.

\begin{algorithm}
\SetAlgoLined
\Fn{\Push{$S$, $x$}}{
    \Se{$S.\NumElem = S.\Len$}{
        \textbf{erro} ``stack overflow''\;
    }
    \InsereVetor{$S$, $x$, $S.\NumElem$}\;
}

\Fn{\Pop{$S$}}{
    \Se{$S.\NumElem = 0$}{
        \textbf{erro} ``underflow''\;
    }
    $x \gets S[S.\NumElem - 1]$\;
    \RemoveVetor{$S$, $S.\NumElem - 1$}\;
    \Retorna{$x$}\;
}
\end{algorithm}

Na implementação por listas encadeadas, o topo da pilha é colocado no início da lista.
Isso porque o procedimento \RemoveLista recebe como argumento o membro que antecede
aquele que será removido, de forma a impossibilitar remoções eficientes.

\begin{algorithm}
\SetAlgoLined
\Fn{\Push{$S$, $x$}}{
    \InsereLista{$S$, $S.\Head$, $x$}\;
}

\Fn{\Pop{$S$}}{
    \Se{$S.\Head.\Prox = \nil$}{
        \textbf{erro} ``underflow''\;
    }
    $x \gets S.\Head.\Prox$\;
    \RemoveLista{$S$, $S.\Head$}\;
    \Retorna{$*x$}\;
}
\end{algorithm}
